% GENERATED FILE. Edit canonical lessons, then run npm run generate:books.
% canonical-source-hash: fnv1a64:35e27209ed1d13e2
% canonical-lessons: FR-C525-zapper, FR-C525-telecommande, FR-C525-emission, FR-C525-documentaire, FR-C525-rediffuser

\chapter{Devant l'écran}
\label{ch:fr-devant-l-ecran}
\hlchaptermodality{525}

\begin{chapteropening}
\textbf{By the end of this chapter:} \emph{I can talk about watching television.}

Complete the last lesson of chapter 525: I can talk about watching television.
\end{chapteropening}

\section[zapper]{zapper — to channel-hop; to skip}

\label{lesson:FR-C525-zapper}



\begin{quote}
Before the new one: say the French for (film) director, then the French for shooting.
\end{quote}

\subsection*{You'll want to know: zapper}
\textbf{zapper} — \textquotedblleft{}to channel-hop; to skip\textquotedblright{}.

\begin{quote}
Le soir, il zappe d'une chaîne à l'autre.
\end{quote}

\begin{quote}
\emph{In the evening he flicks from one channel to another.}
\end{quote}

Informally, to forget or skip: \emph{j'ai zappé le rendez-vous}.

\begin{cousinweb}
\emph{zapper} was borrowed from English \emph{zap} (EtymDB), an imitative word.
\end{cousinweb}

\begin{grammarlens}[title={What you've built}]
One more word for telling, explaining and describing at B1.
\end{grammarlens}

\subsection*{Your turn}
\begin{itemize}
\raggedright
  \item \emph{Say it:} \emph{zapper}
  \item \emph{Say it:} \emph{zapper}, once more
  \item \emph{Say it:} a sentence of your own with \emph{zapper}
  \item \emph{Recall:} say the French for (film) director, then the French for shooting, then say \emph{zapper} again
\end{itemize}

\begin{tcolorbox}[breakable,colback=teal!4,colframe=teal!35!black,title={Before you move on}]
What does zapper mean? (\textquotedblleft{}To channel-hop; to skip\textquotedblright{}.) Say it once more.
\end{tcolorbox}
\section[la télécommande]{la télécommande — remote control}

\label{lesson:FR-C525-telecommande}



\begin{quote}
Before the new one: say the French for shooting, then the French for to channel-hop.
\end{quote}

\subsection*{You'll want to know: la télécommande}
\textbf{la télécommande} — \textquotedblleft{}remote control\textquotedblright{}.

\begin{quote}
Personne ne trouve la télécommande.
\end{quote}

\begin{quote}
\emph{Nobody can find the remote control.}
\end{quote}

Informally also \emph{la zapette}.

\begin{cousinweb}
\emph{la télécommande} is built from Greek \emph{tēle} \textquotedblleft{}far\textquotedblright{} and \emph{commande} \textquotedblleft{}command\textquotedblright{}, which goes back, through \emph{commander}, to Latin \emph{commendāre} \textquotedblleft{}to entrust\textquotedblright{}, built on \emph{mandāre} \textquotedblleft{}to entrust, to order\textquotedblright{} (Lewis \& Short).
\end{cousinweb}

\begin{grammarlens}[title={What you've built}]
One more word for telling, explaining and describing at B1.
\end{grammarlens}

\subsection*{Your turn}
\begin{itemize}
\raggedright
  \item \emph{Say it:} \emph{la télécommande}
  \item \emph{Say it:} \emph{la télécommande}, once more
  \item \emph{Say it:} a sentence of your own with \emph{la télécommande}
  \item \emph{Recall:} say the French for shooting, then the French for to channel-hop, then say \emph{la télécommande} again
\end{itemize}

\begin{tcolorbox}[breakable,colback=teal!4,colframe=teal!35!black,title={Before you move on}]
What does la télécommande mean? (\textquotedblleft{}Remote control\textquotedblright{}.) Say it once more.
\end{tcolorbox}
\section[l'émission]{l'émission — (television or radio) programme; emission}

\label{lesson:FR-C525-emission}



\begin{quote}
Before the new one: say the French for to channel-hop, then the French for remote control.
\end{quote}

\subsection*{You'll want to know: l'émission}
\textbf{l'émission} — \textquotedblleft{}(television or radio) programme; emission\textquotedblright{}.

\begin{quote}
Cette émission de cuisine passe tous les jeudis.
\end{quote}

\begin{quote}
\emph{This cooking programme is on every Thursday.}
\end{quote}

Feminine. Also emissions: \emph{les émissions de CO2}.

\begin{cousinweb}
\emph{l'émission} is a learned borrowing from Latin \emph{ēmissiō} \textquotedblleft{}a sending out\textquotedblright{} (EtymDB), from \emph{ēmittere}, built on \emph{mittere} \textquotedblleft{}to send\textquotedblright{} (Lewis \& Short).
\end{cousinweb}

\begin{grammarlens}[title={What you've built}]
One more word for telling, explaining and describing at B1.
\end{grammarlens}

\subsection*{Your turn}
\begin{itemize}
\raggedright
  \item \emph{Say it:} \emph{l'émission}
  \item \emph{Say it:} \emph{l'émission}, once more
  \item \emph{Say it:} a sentence of your own with \emph{l'émission}
  \item \emph{Recall:} say the French for to channel-hop, then the French for remote control, then say \emph{l'émission} again
\end{itemize}

\begin{tcolorbox}[breakable,colback=teal!4,colframe=teal!35!black,title={Before you move on}]
What does l'émission mean? (\textquotedblleft{}(television or radio) programme; emission\textquotedblright{}.) Say it once more.
\end{tcolorbox}
\section[le documentaire]{le documentaire — documentary}

\label{lesson:FR-C525-documentaire}



\begin{quote}
Before the new one: say the French for remote control, then the French for (television or radio) programme.
\end{quote}

\subsection*{You'll want to know: le documentaire}
\textbf{le documentaire} — \textquotedblleft{}documentary\textquotedblright{}.

\begin{quote}
Hier, nous avons vu un documentaire sur les baleines.
\end{quote}

\begin{quote}
\emph{Yesterday we watched a documentary about whales.}
\end{quote}

Also an adjective: \emph{un film documentaire}.

\begin{cousinweb}
\emph{le documentaire} is built on \emph{document} (EtymDB), a learned borrowing from Latin \emph{documentum} \textquotedblleft{}a lesson, a proof\textquotedblright{}, from \emph{docēre} \textquotedblleft{}to teach\textquotedblright{} (Lewis \& Short).
\end{cousinweb}

\begin{grammarlens}[title={What you've built}]
One more word for telling, explaining and describing at B1.
\end{grammarlens}

\subsection*{Your turn}
\begin{itemize}
\raggedright
  \item \emph{Say it:} \emph{le documentaire}
  \item \emph{Say it:} \emph{le documentaire}, once more
  \item \emph{Say it:} a sentence of your own with \emph{le documentaire}
  \item \emph{Recall:} say the French for remote control, then the French for (television or radio) programme, then say \emph{le documentaire} again
\end{itemize}

\begin{tcolorbox}[breakable,colback=teal!4,colframe=teal!35!black,title={Before you move on}]
What does le documentaire mean? (\textquotedblleft{}Documentary\textquotedblright{}.) Say it once more.
\end{tcolorbox}
\section[rediffuser]{rediffuser — to repeat, to rebroadcast}

\label{lesson:FR-C525-rediffuser}



\begin{quote}
Before the new one: say the French for (television or radio) programme, then the French for documentary.
\end{quote}

\subsection*{You'll want to know: rediffuser}
\textbf{rediffuser} — \textquotedblleft{}to repeat, to rebroadcast\textquotedblright{}.

\begin{quote}
La chaîne rediffuse ce vieux film chaque été.
\end{quote}

\begin{quote}
\emph{The channel repeats this old film every summer.}
\end{quote}

\emph{Une rediffusion} — a repeat.

\begin{cousinweb}
\emph{rediffuser} is built from \emph{re-} and \emph{diffuser} \textquotedblleft{}to broadcast\textquotedblright{}, from Latin \emph{diffundere} \textquotedblleft{}to pour out, to spread\textquotedblright{}, built on \emph{fundere} \textquotedblleft{}to pour\textquotedblright{} (Lewis \& Short).
\end{cousinweb}

\begin{grammarlens}[title={What you've built}]
One more word for telling, explaining and describing at B1.
\end{grammarlens}

\subsection*{Your turn}
\begin{itemize}
\raggedright
  \item \emph{Say it:} \emph{rediffuser}
  \item \emph{Say it:} \emph{rediffuser}, once more
  \item \emph{Say it:} a sentence of your own with \emph{rediffuser}
  \item \emph{Recall:} say the French for (television or radio) programme, then the French for documentary, then say \emph{rediffuser} again
\end{itemize}

\begin{tcolorbox}[breakable,colback=teal!4,colframe=teal!35!black,title={Before you move on}]
What does rediffuser mean? (\textquotedblleft{}To repeat, to rebroadcast\textquotedblright{}.) Say it once more.
\end{tcolorbox}
